\section{The two index rules, and what the documents do not settle}
\label{sec:rules}

\subsection{The optimised rule}
The minimum-volatility index is the solution to a quadratic programme. With
$\wvec$ the index weights, $\bvec$ the parent index weights, $\Sig$ the
provider's estimated asset covariance matrix, $S$ a sector-membership matrix
and $\Xsty$ the standardised style exposures, the published constraint set is

\begin{align}
 \min_{\wvec}\ & \wvec^{\top}\Sig\wvec \label{eq:qp}\\
 \text{s.t.}\ & \ind^{\top}\wvec = 1, \nonumber\\
 & 0 \le \wvec_i \le \min(1.5\%,\ 20\,\bvec_i) \quad \forall i, \nonumber\\
 & \wvec_i = 0 \ \text{ or } \ \wvec_i \ge 0.05\% \quad \forall i, \nonumber\\
 & \bigl|S^{\top}(\wvec-\bvec)\bigr| \le 5\% \quad \text{per sector},
   \nonumber\\
 & \bigl|\Xsty_{-v}{}^{\!\top}(\wvec-\bvec)\bigr| \le 0.25 \quad
   \text{per non-exempt style}, \nonumber\\
 & \tfrac{1}{2}\lVert \wvec-\wvec_{\rm prev}\rVert_1 \le 10\%. \nonumber
\end{align}

The index is reviewed semi-annually, effective the last business day of May and
November, and the covariance matrix is dated to the end of the month preceding
the review rather than to the review itself \citep{msci2012minvol}. Where the
constraint set admits no feasible point the methodology prescribes a relaxation
sequence, loosening turnover first, then the minimum holding, then the weight
cap and sector band together, and never the style bands; Appendix~\ref{app:ladder}
gives the sequence as implemented.

Three features of \eqref{eq:qp} deserve note before the replication. The
semi-continuous condition on $\wvec_i$ makes the programme mixed-integer: a
holding is either absent or at least five basis points, with nothing in
between, which rules out the continuum of dust positions a plain quadratic
programme would return. The turnover constraint is written against
$\wvec_{\rm prev}$, which is the previous book after it has drifted with
prices, not the previous optimisation target. And $\Xsty_{-v}$ carries a
subscript doing a great deal of work: it is the style block with the volatility
factor removed, and that removal is the whole mechanism by which the programme
is allowed to produce a low-volatility portfolio rather than merely a
constrained one.

\subsection{The heuristic rule}
The low-volatility index is specified in two sentences. From the parent
constituents, select the hundred securities with the lowest realised volatility
over the trailing year of trading days; weight each in proportion to the
reciprocal of that volatility, normalised to sum to one
\citep{invesco2024splv497k}. Formally, with $\sigma_i$ the trailing volatility
and $\mathcal{L}$ the hundred lowest,

\begin{equation}
 \wvec_i = \frac{\sigma_i^{-1}}{\sum_{j\in\mathcal{L}}\sigma_j^{-1}}
 \quad \text{for } i\in\mathcal{L}, \qquad \wvec_i = 0 \text{ otherwise.}
 \label{eq:splv}
\end{equation}

There is no covariance matrix, no sector constraint, no style constraint and no
turnover limit. Rebalancing is quarterly. The rule uses only the diagonal of
the covariance matrix, and only its square root at that.

% Table 1. Descriptive, not empirical: every cell is a provision read off the
% primary documents cited in the note, so there is no frozen output behind it
% and make_tables.py does not generate it. Values must be changed only against
% the source documents.
\begin{table}[tbp]
\caption{The two published rules, side by side.}
\label{tab:rules}
\centering\small
\setlength{\tabcolsep}{5pt}
\begin{tabularx}{\textwidth}{@{}lYY@{}}
\toprule
 & Optimised index & Heuristic index\\
\midrule
Objective        & Minimise predicted variance & None stated\\
Risk input       & Full covariance matrix from a commercial factor model
                 & Trailing one-year realised volatility, per stock\\
Selection        & Implied by the programme & The 100 lowest-volatility
                   parent members\\
Weighting        & Implied by the programme & Proportional to
                   $1/\sigma_i$\\
\addlinespace
Weight cap       & $\min(1.5\%,\ 20\times$ parent weight$)$ & None\\
Minimum holding  & 0.05\%, or zero & None\\
Sector band      & $\pm5\%$ active per sector & None\\
Style band       & $\pm0.25$ SD active per style, volatility exempt & None\\
Turnover cap     & 10\% one-way per review & None\\
\addlinespace
Review frequency & Semi-annual & Quarterly\\
Effective date   & Last business day of May and November
                 & Third Friday of the review month\\
Risk-model date  & End of the month preceding the review & Not applicable\\
Infeasibility    & Documented relaxation sequence & Cannot arise\\
\bottomrule
\end{tabularx}
\notetext{Provisions for the optimised index are from its provider's published
methodology \citep{msci2012minvol}; those for the heuristic index are from the
tracking fund's summary prospectus \citep{invesco2024splv497k}, which states the
selection and weighting rules in the issuer's own words. The heuristic index's
review calendar and effective date are the one entry in this table taken from
published convention rather than from a methodology document: the provider's
methodology PDF refused the request from the run environment and was not
retrieved, which Section~\ref{sec:limits} records as an unverified parameter.
The style band is the row that does not survive translation into another risk
model, and Section~\ref{sec:rules} takes it up.}
\end{table}


\subsection{What the style band does not settle}
Every constraint in \eqref{eq:qp} except one translates without ambiguity into
a replicator's own model. A 1.5\% cap is a 1.5\% cap. The sector band needs a
sector classification, and using a coarser one than the provider is a
documented approximation with a known direction of error. The turnover cap and
the minimum holding are numbers.

The style band is different in kind. It is stated in units of standardised
exposure to factors the methodology names but does not define, and the name it
gives the exempt factor is the name of a component of a proprietary model. The
document governing the index examined here exempts \emph{the Barra Volatility
risk index factor}, in the singular. A public-data factor model built to the
conventional recipe has two descriptors with an equally good claim to that
name: the beta of a stock's returns to the parent index, and the volatility of
the residual from that same regression. The first is systematic risk, the
second is the part orthogonal to it. Both are volatility in any ordinary sense,
and the published text does not choose between them.

The choice is not resolved by looking elsewhere in the provider's
documentation, though it is tempting to try. A later methodology covering the
provider's \emph{advanced} minimum-volatility family does enumerate the
exemption as two named factors, beta and residual volatility
\citep{msci2025minvoladvanced}. That document governs a different index family
from the one tracked here, so it cannot be read as a statement about this
index's constraint set. It establishes that the provider's own vocabulary
admits both readings, which is the opposite of settling the question.

A replicator therefore faces a genuine fork, and both branches are defensible.
Exempting residual volatility alone is the literal reading: one named factor,
one exemption, and residual volatility is the descriptor a risk model would
most naturally call volatility once beta has been accounted for separately.
Exempting beta as well is the functional reading: the purpose of the exemption
is to let the optimiser reduce risk, the optimiser's principal lever for
reducing risk is to hold low-beta stocks, and a $\pm0.25$ standard deviation
band on beta forecloses most of that lever before the programme starts.

We implement both and report both. Section~\ref{sec:repro} shows that the fork
is worth 2.10 percentage points of tracking error, and
Section~\ref{sec:estimator} that the choice among four covariance estimators is
worth a fortieth of that. The point is not that one reading is correct. It is
that a reader of the methodology document cannot tell which the index uses, and
that this single unresolved sentence dominates every estimation decision made
downstream of it.
